\documentclass[11pt,a4paper]{article}

\usepackage[margin=2.5cm]{geometry}
\usepackage{amsmath,amssymb}
\usepackage{listings}
\usepackage{xcolor}
\usepackage{hyperref}
\usepackage{booktabs}
\usepackage{microtype}
\usepackage{parskip}

\hypersetup{
  colorlinks=true,
  linkcolor=black,
  urlcolor=black,
  citecolor=black,
  pdftitle={Vero: A Coherence-Guaranteed Language for Neural Architecture Specification},
  pdfauthor={T. Billy Ekkerd}
}

% Vero code listings
\definecolor{veroblue}{RGB}{60,80,140}
\definecolor{verogrey}{RGB}{120,120,130}
\definecolor{verodark}{RGB}{20,20,30}

\lstdefinelanguage{Vero}{
  keywords={model,nodes,edges,phases,boundaries,weight,genome,rule,true,false,provable,unprovable,coherent,dissonant,now,substrate},
  keywordstyle=\color{veroblue}\bfseries,
  comment=[l]{\#},
  comment=[l]{//},
  commentstyle=\color{verogrey}\itshape,
  stringstyle=\color{veroblue},
  basicstyle=\ttfamily\small,
  showstringspaces=false,
  breaklines=true,
  frame=single,
  rulecolor=\color{verogrey!40},
  backgroundcolor=\color{gray!5},
  numbers=left,
  numberstyle=\tiny\color{verogrey},
  numbersep=6pt,
}

\title{\textbf{Vero: A Coherence-Guaranteed Language\\for Neural Architecture Specification}}

\author{
  T.\ Billy Ekkerd \\
  USC Labs --- Unified Systems Chain \\
  Johannesburg, South Africa \\
  \texttt{admin@unifiedsystemschain.com} \\
  ORCID: 0009-0002-5787-3980
}

\date{April 2026}

\begin{document}

\maketitle

\begin{abstract}
We present \textbf{Vero}, a declarative language for specifying neural network architectures
with compile-time coherence guarantees. Unlike existing configuration formats such as JSON or
YAML, which carry no structural guarantees, a Vero program cannot be compiled if it violates
coherence. Every declaration in Vero is annotated with four base properties --- \emph{truth},
\emph{harmony}, \emph{time}, and \emph{place} --- which are verified statically before the
program is lowered to an executable intermediate representation. Dissonant declarations do not
produce warnings; they produce compile errors. We describe the language design, its type system,
the coherence rules, and the current implementation of the Vero Lite compiler, which lowers
Vero source to a JSON intermediate representation consumed by downstream neural runtimes.
The full compiler, specification, and examples are released as open source under the
MIT/Apache-2.0 dual license at \url{https://github.com/U-S-C-Labs/Vero}.
\end{abstract}

% ─────────────────────────────────────────────────────────────────────────────
\section{Introduction}
% ─────────────────────────────────────────────────────────────────────────────

Neural architecture specification is a largely unsolved problem. Frameworks such as
PyTorch~\cite{paszke2019pytorch} and JAX~\cite{bradbury2018jax} describe architectures
imperatively in Python. Configuration formats such as JSON, YAML, and TOML describe
them declaratively but carry no guarantees --- a configuration file that specifies an
incoherent architecture is syntactically identical to a valid one until runtime.

The consequences are practical: silent misconfiguration, weight bounds violations that
propagate through training, and architecture descriptions that cannot be reasoned about
formally. There is no notion of a \emph{wrong} weight annotation in these formats, because
the formats have no semantics beyond key-value storage.

Vero takes a different position: \textbf{a neural architecture specification that cannot be
verified is not a specification}. It is data.

The language draws its design philosophy from three observations:

\begin{enumerate}
  \item Every neural parameter has a \emph{location} (where in the substrate it operates),
        a \emph{moment} (the phase of computation in which it is valid), a \emph{truth value}
        (whether its value can be proven consistent with its context), and a \emph{harmony
        state} (whether it coheres structurally with the system around it).

  \item These four properties --- truth, harmony, time, place --- are not metadata.
        They are \emph{types}. The compiler's primary obligation is to verify them.

  \item A declaration that is structurally dissonant with its context should not compile.
        Not as a warning. As an impossibility.
\end{enumerate}

This paper describes the design and implementation of Vero Lite, the current release of the
language. Section~\ref{sec:design} describes the language design and four-property type system.
Section~\ref{sec:grammar} gives the grammar and worked examples. Section~\ref{sec:coherence}
defines the coherence rules enforced at compile time. Section~\ref{sec:ir} describes the
intermediate representation. Section~\ref{sec:implementation} covers the implementation.
Section~\ref{sec:related} situates Vero relative to related work.
Section~\ref{sec:conclusion} concludes.

% ─────────────────────────────────────────────────────────────────────────────
\section{Language Design}
\label{sec:design}
% ─────────────────────────────────────────────────────────────────────────────

\subsection{The Four Base Properties}

Every declaration in Vero carries four irreducible properties:

\begin{center}
\begin{tabular}{lll}
\toprule
\textbf{Property} & \textbf{Values} & \textbf{Meaning} \\
\midrule
\texttt{truth}   & \texttt{provable} $\mid$ \texttt{unprovable} & Is the value consistent with its context? \\
\texttt{harmony} & \texttt{coherent} $\mid$ \texttt{dissonant}  & Does the declaration cohere structurally? \\
\texttt{time}    & \texttt{now}                                  & The computational phase of validity \\
\texttt{place}   & \texttt{substrate\_current}                   & The substrate location of operation \\
\bottomrule
\end{tabular}
\end{center}

These are not optional annotations. A Vero declaration without all four properties cannot
be well-typed. The compiler verifies each before lowering to the IR.

\subsection{Coherence as a Type Constraint}

In conventional type systems, a type error means a value has the wrong shape.
In Vero, a coherence error means a declaration is \emph{structurally inconsistent with its
context} --- independent of whether its values are numerically valid.

The key coherence contracts are:

\begin{itemize}
  \item \textbf{Truth-Harmony contract}: a declaration with \texttt{truth: unprovable} cannot
        simultaneously claim \texttt{harmony: coherent}. An unprovable value cannot assert
        structural coherence.
  \item \textbf{Dissonance rejection}: \texttt{harmony: dissonant} does not compile, in any
        context. Dissonance is not a valid operational state.
  \item \textbf{Bounds contract}: the declared \texttt{value} must satisfy
        $\texttt{lower} \leq \texttt{value} \leq \texttt{upper}$. A value that violates its
        own declared bounds is incoherent by definition.
  \item \textbf{Identity contract}: every declaration must have a non-empty, non-default name.
        Anonymous declarations cannot be reasoned about.
\end{itemize}

These rules are checked exhaustively --- all violations are reported, not just the first.

\subsection{Design Principle: Boundary Honesty}

Vero encodes a principle we term \emph{boundary honesty}: a measurement or computation that
falls outside its declared domain must not return an incorrect finite value. It must return
infinity or report the measurement as unknown.

This is formalised in the \texttt{boundaries} block, which declares explicit lower and upper
bounds for runtime measurements. Any measurement outside these bounds is reported as
infinity, not clamped or silently discarded. This is a direct expression of Law 17 of the
Universal Laws specification~\cite{ekkerd2026laws} encoded in the language semantics:
\emph{an instrument that reports infinity at its boundary is not wrong --- it has reached its edge.}

% ─────────────────────────────────────────────────────────────────────────────
\section{Grammar and Examples}
\label{sec:grammar}
% ─────────────────────────────────────────────────────────────────────────────

\subsection{Block Structure}

Vero uses a block-structured syntax. A program is a sequence of top-level blocks:

\begin{center}
\begin{tabular}{ll}
\toprule
\textbf{Block} & \textbf{Purpose} \\
\midrule
\texttt{model \{ \}}      & Architecture hyperparameters \\
\texttt{nodes \{ \}}      & Node-level overrides \\
\texttt{edges \{ \}}      & Edge-level overrides \\
\texttt{phases \{ \}}     & Phase scheduling constraints \\
\texttt{boundaries \{ \}} & Runtime measurement domains \\
\texttt{weight \{ \}}     & Named, annotated weight declarations \\
\texttt{genome \{ \}}     & Evolutionary search configuration \\
\texttt{rule \{ \}}       & Symbolic plasticity rules \\
\bottomrule
\end{tabular}
\end{center}

Comments are introduced with \texttt{\#} or \texttt{//} and extend to end of line.
Key-value lines take the form \texttt{key: value}.

\subsection{Minimal Example}

\begin{lstlisting}[language=Vero, caption={Minimal Vero program specifying a spiking network configuration.}]
# Vero-lite: minimal character-level model

model char_mvp {
  hidden_size:      256
  ticks_per_char:   6
  p_spike:          0.35
  pulse_current:    3.0
  plasticity_every: 64
  lr:               0.05
}

phases {
  inference_topology_locked:   true
  plasticity_topology_mutable: true
}

boundaries {
  membrane_sum_lower: 0
  membrane_sum_upper: 5000000
  loss_lower:         0
  loss_upper:         50000
}
\end{lstlisting}

\subsection{Weight Declaration with Coherence Annotations}

\begin{lstlisting}[language=Vero, caption={Weight declaration with full coherence annotation. A value outside declared bounds, or an unprovable value claiming coherent harmony, will not compile.}]
weight plasticity_lr {
  value:    0.07
  verified: bounds(0.0, 1.0)
  truth:    provable
  harmony:  coherent
  time:     now
  place:    substrate.current
  compassion: harmony
}
\end{lstlisting}

\subsection{Genome Declaration}

\begin{lstlisting}[language=Vero, caption={Genome block specifying evolutionary search parameters.}]
genome {
  population_size:       50
  speciation_threshold:  3.0
  survival_rate:         0.2
  crossover_rate:        0.75
  weight_perturb_rate:   0.8
  weight_perturb_sigma:  0.1
  add_node_prob:         0.03
  add_edge_prob:         0.05
  stagnation_limit:      15
  elite_per_species:     1
  seed:                  0
  truth:    provable
  harmony:  coherent
}
\end{lstlisting}

\subsection{B-Spline Edge Configuration}

KAN edges~\cite{liu2024kan} use B-spline basis functions. The number of control points per
edge is $n_\text{ctrl} = G + k$, where $G$ is the number of interior grid intervals
(\texttt{kan\_n\_grid}) and $k$ is the polynomial degree (\texttt{kan\_degree}):

\begin{center}
\begin{tabular}{cccl}
\toprule
$G$ & $k$ & $n_\text{ctrl}$ & Notes \\
\midrule
5 & 3 & 8  & Default (cubic) \\
8 & 3 & 11 & Higher resolution \\
1 & 1 & 2  & Minimal (linear) \\
7 & 2 & 9  & Quadratic \\
\bottomrule
\end{tabular}
\end{center}

% ─────────────────────────────────────────────────────────────────────────────
\section{Coherence Rules}
\label{sec:coherence}
% ─────────────────────────────────────────────────────────────────────────────

The coherence checker runs over the full program before any IR is emitted.
It returns \emph{all} violations, not just the first. A program with any coherence
violation produces no output.

\subsection{Weight Coherence}

For each \texttt{WeightDecl}:

\begin{enumerate}
  \item $\texttt{lower} \leq \texttt{value} \leq \texttt{upper}$
  \item $\neg(\texttt{truth} = \texttt{unprovable} \wedge \texttt{harmony} = \texttt{coherent})$
  \item $\texttt{harmony} \neq \texttt{dissonant}$
  \item $\texttt{name} \neq \varnothing \wedge \texttt{name} \neq \texttt{"unnamed"}$
\end{enumerate}

\subsection{Genome Coherence}

For each \texttt{GenomeDecl}:

\begin{enumerate}
  \item $\texttt{harmony} \neq \texttt{dissonant}$
  \item $\texttt{population\_size} > 0$
  \item $0 < \texttt{survival\_rate} \leq 1.0$
\end{enumerate}

\subsection{Rule Coherence}

For each \texttt{RuleDecl}:

\begin{enumerate}
  \item $\texttt{name} \neq \varnothing \wedge \texttt{name} \neq \texttt{"unnamed"}$
  \item $\texttt{harmony} \neq \texttt{dissonant}$
  \item $\texttt{expression} \neq \varnothing$
\end{enumerate}

% ─────────────────────────────────────────────────────────────────────────────
\section{Intermediate Representation}
\label{sec:ir}
% ─────────────────────────────────────────────────────────────────────────────

Vero Lite compiles to a JSON intermediate representation (IR). The top-level structure is:

\begin{lstlisting}[language={},caption={Vero IR top-level structure (JSON).}]
{
  "version": 1,
  "model":      { "kind": "char_mvp", ... },
  "phases":     { "inference_topology_locked": true, ... },
  "boundaries": { "membrane_sum": {...}, "loss": {...} },
  "weights":    [ { "name": "...", "value": ..., ... } ],
  "genome":     { ... } | null,
  "rules":      [ { "name": "...", "expression": "..." } ]
}
\end{lstlisting}

The IR is consumed by runtime backends. The current backend targets CPU execution via a
spiking neural network runtime implemented in Rust. A GPU backend (WebGPU/wgpu) is under
development. The IR design intentionally does not embed backend-specific information ---
backend selection and lowering are downstream concerns, not language concerns.

The IR carries all four Vero base annotations (\texttt{truth}, \texttt{harmony},
\texttt{time}, \texttt{place}) on each \texttt{WeightDecl} and \texttt{GenomeDecl}.
These annotations are preserved through lowering so that runtime systems can inspect
and enforce them independently of the compiler.

% ─────────────────────────────────────────────────────────────────────────────
\section{Implementation}
\label{sec:implementation}
% ─────────────────────────────────────────────────────────────────────────────

The Vero Lite compiler is implemented in Rust (edition 2024) as two crates:

\begin{itemize}
  \item \textbf{\texttt{vero\_ir}} --- the intermediate representation types and
        coherence checker (\textasciitilde{}250 lines).
  \item \textbf{\texttt{vero}} --- the parser, compiler CLI, and lowering pass
        (\textasciitilde{}450 lines).
\end{itemize}

The compiler binary is invoked as:

\begin{lstlisting}[language=bash]
vero compile <input.vero> --out <output.json>
\end{lstlisting}

The parser is a single-pass line-oriented parser with no external parser generator
dependency. This is a deliberate design choice: the grammar is simple enough that a
hand-written parser is easier to audit than a generated one, and auditability is a
first-class concern for a language intended for safety-critical neural governance.

The coherence checker uses Rust's type system extensively --- the \texttt{TruthAnnotation}
and \texttt{HarmonyAnnotation} enums make invalid combinations representable but
statically detectable, and the checker exhaustively matches against them. All violations
are collected into a \texttt{Vec<CoherenceError>} and returned together, so authors see
all problems at once rather than fixing them one at a time.

The full implementation, specification, and examples are available at:\\
\url{https://github.com/U-S-C-Labs/Vero}

% ─────────────────────────────────────────────────────────────────────────────
\section{Related Work}
\label{sec:related}
% ─────────────────────────────────────────────────────────────────────────────

\textbf{Neural architecture description formats.}
ONNX~\cite{bai2019onnx} provides a portable format for neural network graphs but
focuses on inference-time exchange rather than specification-time verification.
It carries no coherence guarantees. Keras configuration files and PyTorch
\texttt{state\_dict} exports are similarly unverified serialisation formats.

\textbf{Formal verification of neural networks.}
A body of work applies formal methods to verify \emph{properties of trained networks}
(robustness, fairness, safety bounds) post-hoc~\cite{katz2017reluplex,liu2021algorithms}.
Vero addresses a different problem: verifying the \emph{specification} of the architecture
before any training or execution occurs.

\textbf{Verified systems programming.}
Verus~\cite{lattuada2023verus} extends Rust with SMT-backed proof obligations for
systems code. Vero draws philosophical inspiration from Verus but operates at a higher
level of abstraction --- it verifies architectural coherence rather than memory safety
or functional correctness of general programs.

\textbf{Domain-specific languages for ML.}
Halide~\cite{ragan2013halide} and TVM~\cite{chen2018tvm} provide DSLs for specifying
and optimising tensor computations. These languages are concerned with computational
efficiency; they have no notion of architectural coherence as a first-class property.

To our knowledge, no existing language treats neural architecture \emph{coherence} ---
the structural consistency of declarations with their declared truth, harmony, temporal,
and spatial context --- as a compile-time type constraint.

% ─────────────────────────────────────────────────────────────────────────────
\section{Conclusion}
\label{sec:conclusion}
% ─────────────────────────────────────────────────────────────────────────────

We have presented Vero, a declarative language for neural architecture specification
with compile-time coherence guarantees. The core contribution is the identification of
four irreducible properties of neural declarations --- truth, harmony, time, and place ---
and their formalisation as compile-time type constraints enforced by the coherence checker
before any IR is emitted.

The current release, Vero Lite, implements the parser, coherence checker, and JSON IR
lowering pass in approximately 700 lines of auditable Rust. It is the language layer of
a larger research programme into coherence-guaranteed neural architectures. Subsequent
work will extend the language with dependent types over the four base properties, formal
proof obligations on weight declarations, and a full \texttt{VeroMind} genome specification
for spatially-embedded neuroevolution.

The design principle underlying all of this work is simple:
\emph{a system that cannot verify its own coherence cannot be trusted.}
Vero is an attempt to make that verification the first thing the compiler does,
not the last.

% ─────────────────────────────────────────────────────────────────────────────
\begin{thebibliography}{99}
% ─────────────────────────────────────────────────────────────────────────────

\bibitem{paszke2019pytorch}
A.\ Paszke et al.,
``PyTorch: An Imperative Style, High-Performance Deep Learning Library,''
\emph{Advances in Neural Information Processing Systems}, 2019.

\bibitem{bradbury2018jax}
J.\ Bradbury et al.,
``JAX: Composable Transformations of Python+NumPy Programs,''
2018. \url{http://github.com/jax-ml/jax}

\bibitem{liu2024kan}
Z.\ Liu et al.,
``KAN: Kolmogorov-Arnold Networks,''
\emph{arXiv preprint arXiv:2404.19756}, 2024.

\bibitem{bai2019onnx}
J.\ Bai et al.,
``ONNX: Open Neural Network Exchange,''
2019. \url{https://github.com/onnx/onnx}

\bibitem{katz2017reluplex}
G.\ Katz et al.,
``Reluplex: An Efficient SMT Solver for Verifying Deep Neural Networks,''
\emph{CAV}, 2017.

\bibitem{liu2021algorithms}
C.\ Liu et al.,
``Algorithms for Verifying Deep Neural Networks,''
\emph{Foundations and Trends in Optimization}, 2021.

\bibitem{lattuada2023verus}
A.\ Lattuada et al.,
``Verus: Verifying Rust Programs Using Linear Ghost Types,''
\emph{OOPSLA}, 2023.

\bibitem{ragan2013halide}
J.\ Ragan-Kelley et al.,
``Halide: A Language and Compiler for Optimizing Parallelism, Locality, and Recomputation in Image Processing Pipelines,''
\emph{PLDI}, 2013.

\bibitem{chen2018tvm}
T.\ Chen et al.,
``TVM: An Automated End-to-End Optimizing Compiler for Deep Learning,''
\emph{OSDI}, 2018.

\bibitem{ekkerd2026laws}
T.\ B.\ Ekkerd,
``The Complete Architecture of Natural Law --- 17 Laws in Three Parts,''
USC Labs Technical Document, 2026.

\end{thebibliography}

\end{document}
